% Chapter 19: Word Building and One-Word Expressions: Growing Your Vocabulary

\chapter{Word Building and One-Word Expressions: Growing Your Vocabulary}

\section{Pedagogical Purpose}
Having explored synonyms, antonyms, and homophones, learners now learn how new words are \textit{built} from smaller parts. Understanding prefixes, suffixes, and word families helps learners guess the meaning of unfamiliar words and expand their vocabulary efficiently. One-word substitutions also make writing more concise and precise.

\begin{learningobjectives}
The learner can:
\begin{itemize}
    \item identify the root word in a word with a prefix or suffix.
    \item add a prefix to change a word's meaning (e.g., un-, re-, dis-).
    \item add a suffix to change a word's form (e.g., -ful, -less, -ly).
    \item group words that belong to the same word family.
    \item use one-word expressions to replace common longer phrases.
\end{itemize}
\end{learningobjectives}

\section{Prerequisite Check}
\begin{quickcheck}
What smaller word can you find inside ``happiness''? (\textit{happy})
\end{quickcheck}

\begin{rememberbox}
You have learned many individual words already. Now you will learn how words are built from smaller pieces, and how one word can sometimes replace a whole phrase.
\end{rememberbox}

\section{Chapter Opener}
Maya is building a tower of blocks, each labelled with a word part, while Ms. Sharma watches.

\textbf{Maya:} ``Look, Ms. Sharma! I have a block that says `kind,' and if I add this block that says `-ness,' it becomes `kindness'!''

\textbf{Rohan:} ``Can I add `un-' to the front of `happy' to make `unhappy'?''

\textbf{Ms. Sharma:} ``Wonderful! You're both building words the same way English does --- by joining small meaningful parts together. Let's explore this word-building game properly.''

\section{Language Experience}
\textbf{Word-Building Blocks}

\begin{itemize}
    \item kind + ness = kindness
    \item happy $\rightarrow$ un + happy = unhappy
    \item care + ful = careful
    \item care + less = careless
    \item play + er = player
\end{itemize}

\section{Look and Notice}
\begin{thinkbox}
\begin{enumerate}
    \item What is the small, meaningful word hiding inside ``kindness''?
    \item What does adding ``un-'' to the front of ``happy'' do to its meaning?
    \item What is the difference between ``careful'' and ``careless''?
    \item Can you find other words that end in ``-er'' and mean ``a person who does something''?
\end{enumerate}
\end{thinkbox}

\section{Discover / Explicit Teaching}
Many English words are built from a \textbf{root word} plus a \textbf{prefix} (added to the front) or a \textbf{suffix} (added to the end).

\begin{table}[h!]
\centering
\begin{tabular}{|l|l|l|l|}
\hline
\textbf{Part} & \textbf{Position} & \textbf{Job} & \textbf{Example} \\
\hline
Root word & --- & Carries the main meaning & kind, care, play \\
Prefix & Front & Changes meaning, often to opposite & un + happy = unhappy \\
Suffix & End & Changes word form or adds meaning & kind + ness = kindness \\
\hline
\end{tabular}
\end{table}

\section{Grammar Word}
\begin{grammarword}{PREFIX}
A \textbf{prefix} is a word part added to the \textbf{beginning} of a root word to change its meaning.

\textit{Examples:}
\begin{itemize}
    \item un + happy = unhappy (not happy)
    \item re + write = rewrite (write again)
    \item dis + agree = disagree (not agree)
\end{itemize}

\textit{Non-example:} ``-ness'' added to ``kind'' (kindness) is not a prefix --- it is added to the end, so it is a suffix.

\textit{Quick Check:} Add ``un-'' to ``kind.'' What is the new word and its meaning? \textit{(unkind --- not kind)}
\end{grammarword}

\begin{grammarword}{SUFFIX}
A \textbf{suffix} is a word part added to the \textbf{end} of a root word to change its form or meaning.

\textit{Examples:}
\begin{itemize}
    \item care + ful = careful (full of care)
    \item care + less = careless (without care)
    \item teach + er = teacher (a person who teaches)
\end{itemize}

\textit{Non-example:} ``re-'' added to ``write'' (rewrite) is not a suffix --- it is added to the beginning, so it is a prefix.

\textit{Quick Check:} Add ``-less'' to ``hope.'' What is the new word and its meaning? \textit{(hopeless --- without hope)}
\end{grammarword}

\section{Core Explanation}

\subsection*{Prefixes often reverse or change meaning}
agree $\rightarrow$ \textbf{dis}agree (opposite). \textit{Check:} Does the prefix flip the root word's meaning to its opposite, or add ``again'' (re-)?

\subsection*{Suffixes often change word form}
help (verb) $\rightarrow$ help\textbf{ful} (adjective, ``full of help''). \textit{Check:} Does the suffix turn the word into a describing word (adjective) or a naming word (noun)?

\subsection*{Word families share a root}
play, player, playful, playground --- all share the root ``play.'' \textit{Check:} Does the word truly share the meaning of the root, not just similar letters?

\subsection*{One-word expressions}
``a person who cannot hear'' $\rightarrow$ \textbf{deaf}; ``a place where books are kept'' $\rightarrow$ \textbf{library}. \textit{Check:} Is there a single word that captures the same meaning as the phrase?

\section{Worked Example}
\begin{examplebox}
\textbf{Task:} Build a new word by adding the correct prefix or suffix, and explain the meaning: ``agree'' + ``not agreeing.''

\textit{Step 1:} Identify the root word --- ``agree.''
\textit{Step 2:} Identify the meaning needed --- ``not agreeing'' suggests an opposite meaning.
\textit{Step 3:} Choose the correct prefix --- ``dis-'' means ``not,'' so ``dis'' + ``agree'' = ``disagree.''

\textbf{Answer:} \textbf{disagree} --- meaning ``to not agree.''
\end{examplebox}

\section{Guided Practice}
\begin{activitybox}{Activity A: Identification}
Underline the root word and circle the prefix or suffix.
\begin{enumerate}
    \item \textbf{un}happy
    \item care\textbf{ful}
\end{enumerate}
\end{activitybox}

\begin{activitybox}{Activity B: Completion}
Build the correct word.
\begin{enumerate}
    \item play + er = \_\_\_
    \item kind + ness = \_\_\_
\end{enumerate}
\end{activitybox}

\section{Independent Practice}
\begin{activitybox}{Independent Practice}
\begin{enumerate}
    \item Add ``un-'' to three root words of your choice and explain each new meaning.
    \item Add ``-ful'' and ``-less'' to the word ``help'' and explain the difference in meaning.
    \item List three words that belong to the ``teach'' word family.
    \item \textit{Reasoning:} Why do you think knowing common prefixes and suffixes can help you guess the meaning of a word you have never seen before?
\end{enumerate}
\end{activitybox}

\section{Grammar Detective}
\begin{detectivebox}
\textbf{Student sentence:}

``The unkindful boy was carelessness to his little sister.''

\textit{Questions:}
\begin{enumerate}
    \item What is wrong?
    \item What should it be?
    \item Why?
\end{enumerate}

\textit{Guidance:} ``unkindful'' is not a real word; it should be ``unkind.'' ``carelessness'' (a noun) is used where the adjective ``careless'' is needed. The correct sentence is: ``The unkind boy was careless to his little sister.''
\end{detectivebox}

\section{Common Confusion}
\begin{warningbox}
Learners sometimes combine two word-building parts that do not belong together (like ``unkindful'') or use a noun form (``carelessness'') where an adjective (``careless'') is needed.

\textit{How to check:} Ask, ``Does this word part actually exist for this root, and does the resulting word fit the sentence as a describing word, a naming word, or an action word?''
\end{warningbox}

\section{Writing Connection}
Write four sentences, each using one word built with a prefix or suffix from this chapter (e.g., unhappy, careful, careless, rewrite). Underline the built word in each sentence.

\section{Summary}
\begin{summarybox}
\begin{itemize}
    \item Words are often built from a \textbf{root word} plus a \textbf{prefix} (front) or \textbf{suffix} (end).
    \item Prefixes like \textbf{un-}, \textbf{re-}, and \textbf{dis-} change meaning, often to the opposite.
    \item Suffixes like \textbf{-ful}, \textbf{-less}, \textbf{-ly}, and \textbf{-er} change a word's form or add meaning.
    \item Words sharing a root form a \textbf{word family}.
    \item A \textbf{one-word expression} can often replace a longer phrase.
\end{itemize}
\end{summarybox}

\section{Key Terms}
\begin{table}[h!]
\centering
\begin{tabular}{|l|l|}
\hline
\textbf{Term} & \textbf{Simple Meaning} \\
\hline
Root word & The main word that carries the core meaning \\
Prefix & A word part added to the front to change meaning \\
Suffix & A word part added to the end to change form or meaning \\
Word family & A group of words sharing the same root \\
One-word expression & A single word that replaces a longer phrase \\
\hline
\end{tabular}
\end{table}

\begin{selfcheck}
I can:
\begin{itemize}
    \item identify a root word within a longer word.
    \item add a prefix to change a word's meaning.
    \item add a suffix to change a word's form.
    \item group words into a word family.
    \item use a one-word expression instead of a longer phrase.
\end{itemize}
\end{selfcheck}